%% ─────────────────────────────────────────────────────────────────────────────
%% Lettre de motivation — Talvin Ackbaraly × MBDA France
%% Stage : détection de risque de collision sur séquences d'images
%% Même charte que le CV (XCharter, un seul accent, filets fins)
%% Compile: latexmk -pdf lettre-motivation-mbda.tex
%% ─────────────────────────────────────────────────────────────────────────────
\documentclass[11pt, a4paper]{extarticle}

\usepackage[T1]{fontenc}
\usepackage[utf8]{inputenc}
\usepackage[french]{babel}
\usepackage[top=1.6cm, bottom=1.6cm, left=2.2cm, right=2.2cm]{geometry}
\usepackage{XCharter}
\usepackage[letterspace=80]{microtype}
\usepackage{ragged2e}
\usepackage{xcolor}
\usepackage[hidelinks]{hyperref}

\hypersetup{pdftitle={Lettre de motivation — Talvin Ackbaraly}, pdfauthor={Talvin Ackbaraly}}

%% ── Colour system: identique au CV ──────────────────────────────────────────
\definecolor{accent}{HTML}{2563A0}
\definecolor{ink}{HTML}{1A1A1A}
\definecolor{soft}{HTML}{6E6E6E}
\definecolor{hair}{HTML}{C9C9C9}

\hypersetup{colorlinks=true, urlcolor=accent, linkcolor=accent}
\AtBeginDocument{\color{ink}}

\pagestyle{empty}
\setlength{\parindent}{0pt}
\setlength{\JustifyingParindent}{0pt}
\setlength{\parskip}{8pt}
\linespread{1.08}
\microtypesetup{expansion=false}

\begin{document}

%% ── En-tête : expéditeur ─────────────────────────────────────────────────────
%% Coordonnées seulement : l'école et les liens sont déjà dans le CV
{\color{accent}\fontsize{22}{26}\selectfont\scshape\textls[40]{Talvin Ackbaraly}}\par\vspace{3pt}
{\setlength{\parskip}{0pt}
94270 Le Kremlin-Bicêtre\\
\href{tel:+33769389008}{\textcolor{ink}{07 69 38 90 08}}\\
\href{mailto:talvin.ackbaraly@icloud.com}{\textcolor{ink}{talvin.ackbaraly@icloud.com}}\par}
\vspace{-4pt}{\color{hair}\rule{\textwidth}{0.4pt}}\par

%% ── Destinataire et date ─────────────────────────────────────────────────────
\vspace{2pt}
%% Bloc calé contre la marge droite, texte aligné à gauche dans le bloc
\hfill\begin{tabular}[t]{@{}l@{}}
  \textbf{MBDA France}\\
  Direction Technique\\
  1 avenue Réaumur\\
  92350 Le Plessis-Robinson
\end{tabular}

Le \today

\vspace{2pt}
\textbf{Objet : candidature au stage de fin d'études « Détection de risque de collision sur
séquences d'images » — 6 mois à partir de février 2027}

\vspace{2pt}
Madame, Monsieur,

\justifying
Étudiant en dernière année du cycle ingénieur à l'EPITA, dans la majeure Image, je vous adresse
ma candidature pour le stage consacré à la détection de risque de collision sur séquences d'images
au sein de la Direction Technique de MBDA France. Je veux terminer ma formation sur un sujet où
l'algorithme de vision a des conséquences très concrètes. Chez MBDA, la fiabilité d'une méthode
compte au moins autant que sa performance moyenne.

Ma majeure m'a donné des bases solides en traitement d'images, classique comme en deep learning,
et je les ai appliquées à des projets réels. Pour mon projet de fin d'études, mené avec la
Bibliothèque nationale de France, je suis en charge de la détection d'objets. J'ai comparé
plusieurs architectures, réglé les hyperparamètres et évalué les modèles sur un jeu de
validation, jusqu'à 94,7\,\% de mAP50 sur 10\,000 images annotées. C'est la démarche que décrit
votre offre : faire l'état de l'art, retenir les méthodes les plus prometteuses et les évaluer
sur des données de référence. J'ai aussi l'habitude de mettre en regard les deux familles
d'approches. J'ai conçu un système de lecture de plaques d'immatriculation fondé sur HOG, la
morphologie mathématique et une forêt aléatoire, et implémenté un U-Net de A à Z en PyTorch.

C'est l'aspect temporel du sujet qui m'intéresse le plus. Dans mon projet de suivi de tumeurs
cérébrales, j'ai travaillé sur le recalage d'IRM (rigide, affine, B-spline) : il s'agit d'estimer
la transformation entre deux acquisitions, ce qui est proche de l'estimation du mouvement
apparent. Pendant l'état de l'art, j'aimerais étudier l'estimation du temps avant contact à
partir du flot optique et de l'expansion apparente des objets. Je la comparerais ensuite aux
approches par réseaux de neurones qui traitent directement des séquences.

Mon parcours en ingénierie logicielle, avec un stage full-stack et DevOps chez Studiocanal puis
une activité en freelance, me servira pour la partie outillage du stage. Je sais écrire du code
Python propre et reproductible, construire des pipelines d'évaluation et prendre en main de
nouveaux outils, comme ceux que vous utilisez pour générer des séquences de test. Je suis
également à l'aise pour rédiger un rapport scientifique, en français comme en anglais
(TOEIC 845), ce qui compte dans l'environnement international de MBDA.

Je suis disponible pour un stage de six mois à partir de février 2027 et serais heureux de vous
présenter plus en détail ma motivation lors d'un entretien.

Je vous prie d'agréer, Madame, Monsieur, l'expression de mes salutations distinguées.

\vspace{10pt}
\hfill Talvin Ackbaraly\hspace*{1cm}

\end{document}
